\documentclass[11pt]{article}

% ---------------------------------------------------------------------------
%  Data-generation report for the Hubbard-GME ML dataset (gme_qml project)
%  Compile:  pdflatex data_generation_report.tex   (run twice for refs/TOC)
% ---------------------------------------------------------------------------

\usepackage[a4paper,margin=2.3cm]{geometry}
\usepackage{amsmath,amssymb}
\usepackage{booktabs}
\usepackage{longtable}
\usepackage{array}
\usepackage{xcolor}
\usepackage{enumitem}
\usepackage[hidelinks]{hyperref}
\usepackage{fancyvrb}

\newcommand{\code}[1]{\texttt{\small #1}}
\newcommand{\CGME}{\ensuremath{C_{\mathrm{GME}}}}
\newcommand{\rdm}{\ensuremath{\Gamma^{(n)}}}

\title{\textbf{Data Generation, Schema, and Auditing}\\[2pt]
       \large Machine-Learning Dataset for Genuine Multipartite Entanglement\\
       in the 1-D Hubbard Model (\code{gme\_qml} project)}
\author{Hubbard-GME project}
\date{\today}

\begin{document}
\maketitle
\tableofcontents
\bigskip

% ===========================================================================
\section{Purpose and scope}
% ===========================================================================
The genuine-multipartite-entanglement measure \CGME{} of an $n$-body reduced
density matrix (RDM) $\rdm$ of the Hubbard model is computed by a convex-roof
optimisation, which is accurate but expensive (seconds to minutes per mixed
RDM, hours per parameter sweep). The \code{gme\_qml} project trains machine-learning
surrogates that predict \CGME{} directly from $\rdm$. This report documents:
\begin{enumerate}[nosep]
  \item how the training data is \emph{generated} (Section~\ref{sec:gen}),
  \item what \emph{type} of data and \emph{every field} that is stored
        (Section~\ref{sec:schema}),
  \item how the dataset is \emph{audited} for quality (Section~\ref{sec:audit}),
  \item the \emph{commands} that produce and check the dataset
        (Section~\ref{sec:commands}),
  \item the \emph{memory} required as a function of $n$ and $L$
        (Section~\ref{sec:memory}).
\end{enumerate}

Two supervised tasks are targeted:
\begin{center}
\begin{tabular}{@{}lll@{}}
\toprule
Task & Type & Target \\
\midrule
Detection  & binary classification & is \CGME{}$>\varepsilon$ ? \\
Estimation & regression            & the value of \CGME{} \\
\bottomrule
\end{tabular}
\end{center}
The ground truth is the verified \textbf{GCG} (Generalised Conjugate Gradient)
convex-roof solver.

% ===========================================================================
\section{Data-generation pipeline}\label{sec:gen}
% ===========================================================================

\subsection{Physics engine (ground truth)}
For a chosen Hamiltonian $(L, N_e, U, t, \text{mode}, \text{sector})$:
\begin{enumerate}[nosep]
  \item Build the Hubbard Hamiltonian and diagonalise it (Sz$=0$ sector by
        default, or the full $N_e$-electron space with \code{-{}-full-hamiltonian}).
  \item For each eigenstate $\psi_m$ and body order $n$, form the $n$-body RDM
        $\rdm[\psi_m]$ (cluster-averaged over contiguous site clusters for $n<L$).
  \item Combine states as needed: ground state ($T=0$) or a Boltzmann mixture
        $\rdm(T)=\sum_m p_m(T)\,\rdm[\psi_m]$ for thermal data.
  \item Label each RDM with \CGME{} via the GCG convex roof.
\end{enumerate}

\subsection{The RDM bank: the reusable primitive}
The object $\{\rdm[\psi_m]\}$ for \emph{all} eigenstates and all $n$ is the
\emph{RDM bank}. It depends \emph{only} on the Hamiltonian and sector --- not on
temperature, solver, or task --- so it is computed \textbf{once} per
$(L,N_e,U,t,\text{mode},\text{sector})$ and cached in
\code{results/rdm\_pkl\_results/rdm\_bank/}. Every later analysis (ground state,
thermal, ML labelling) reuses it, avoiding re-diagonalisation. To keep the peak
memory independent of the number of worker processes, the eigenvectors are
shared across workers by fork copy-on-write rather than pickled per worker.

\subsection{Two ways to obtain labelled records}
\paragraph{(A) Native generator \code{ml\_generate\_ml\_dataset.py}.}
Sweeps $(L,\,U/t \text{ or } t,\,\text{state},\,\text{cluster})$, diagonalises
once per configuration, and emits one record per RDM. Sampling of $U/t$ is
\emph{stratified}: $\sim$25\,\% weak coupling (log-spaced), $\sim$50\,\% in the
crossover $U/t\in[3,10]$ (linear-dense), $\sim$25\,\% strong coupling
(log-spaced), so the physically interesting crossover is oversampled. This
generator labels with the \textbf{Stiefel} solver and stores rich state features
(spin, energy). It is kept for reference but is \emph{not} the GCG dataset.

\paragraph{(B) GCG builder \code{ml\_ingest\_pkls.py}.}
Produces the \textbf{GCG-only} dataset with three sources:
\begin{center}
\begin{tabular}{@{}lll@{}}
\toprule
\code{-{}-source} & RDM $\rho$ from & Label \CGME{} \\
\midrule
\code{gme}     & linked \code{rdm\_cache/} RDM       & existing \code{C\_gcg} in the GME PKL \\
\code{surface} & reconstructed from the bank + $p_m(T)$ & existing \code{C\_surface}$[n][t,T]$ \\
\code{banks}   & bank RDM ($T{=}0$ GS, or $T{>}0$ mixed) & \textbf{GCG run now} on that $\rho$ \\
\bottomrule
\end{tabular}
\end{center}
\code{gme} and \code{surface} \emph{convert} already-labelled results; \code{banks}
\emph{generates} new labels by running GCG on the banks' RDMs and is the way to
scale the dataset across $L$, $t$, $U$, and $T$. Only the GCG value is used ---
Stiefel results are never ingested and no Stiefel-named field is written.

\subsection{Storage layout}
General, reusable artifacts live under \code{results/rdm\_pkl\_results/}; the ML
dataset lives under \code{gme\_qml/data\_for\_ML/}.
\begin{Verbatim}[fontsize=\small]
results/rdm_pkl_results/
  rdm_bank/     rdmbank_..._{sz0|full}_...pkl   general RDM banks (per Hamiltonian)
  gme_runs/     standalone T=0 GME solver results (bounds + per-k/per-restart logs)
  thermal_runs/ thermal / surface results + per-grid-point checkpoints
  rdm_cache/    single RDMs; state_cache/ diagonalised states

gme_qml/data_for_ML/
  records/      native (Stiefel) records + manifest.csv
  ingested/     GCG dataset: records/*.npz + manifest.csv   <-- used for GCG training
\end{Verbatim}

% ===========================================================================
\section{What is stored: record schema}\label{sec:schema}
% ===========================================================================
Each record is one \code{.npz} file (arrays + config) plus one row in
\code{manifest.csv} (scalar summary for fast filtering). The complete field list
is in \code{data\_dictionary.csv}; the tables below group every field by
category. Column \emph{Dataset} indicates where the field appears: \emph{native}
(Stiefel), \emph{gcg} (ingested), or \emph{both}.

\paragraph{Record file (\code{.npz}) arrays.}
\code{rho\_real}, \code{rho\_imag} ($2^n\times2^n$, \code{float32}) --- the RDM;
\code{eigvals}; \code{config\_json}; and the scalar labels listed under
``GME label'' below. The GCG records add \code{source}, \code{solver},
\code{T}, \code{thermal}.

\subsection{System / Hamiltonian fields}
\begin{longtable}{@{}p{2.4cm}p{1.6cm}p{9.3cm}@{}}
\toprule
Field & Dataset & Description \\ \midrule \endhead
\code{L}            & both   & Number of lattice sites \\
\code{t}            & both   & Hopping amplitude \\
\code{U}            & both   & On-site Coulomb interaction \\
\code{U\_over\_t}   & both   & Coupling ratio $U/t$ \\
\code{mode}         & both   & Topology (uniform/open/nnn/ssh/ladder) \\
\code{n\_electrons} & both   & Number of electrons $N_e$ ($N_e{=}L$ is half-filling) \\
\code{n\_so}        & native & Spin orbitals $=2L$ \\
\code{sector}       & gcg    & Hilbert-space sector: \code{sz0} ($S_z{=}0$) or \code{full} \\
\code{energy}       & native & Energy of the source eigenstate \\
\code{gap}          & native & Spectral gap $E_1-E_0$ \\
\bottomrule
\end{longtable}

\subsection{State / spin observables (native)}
\begin{longtable}{@{}p{2.4cm}p{1.6cm}p{9.3cm}@{}}
\toprule
Field & Dataset & Description \\ \midrule \endhead
\code{state\_idx}      & both   & Eigenstate index (0 = ground); \code{"T\{val\}"} for thermal \\
\code{S2}              & native & $\langle S^2\rangle$ \\
\code{S}               & native & Total spin $S$ \\
\code{multiplicity}    & native & $2S+1$ \\
\code{spin\_label}     & native & Singlet / Doublet / Triplet / \dots \\
\code{Sz}              & native & $z$-projection of spin \\
\code{d\_occ}          & native & Total double occupancy $\langle\sum_j n_{j\uparrow}n_{j\downarrow}\rangle$ \\
\code{d\_occ\_per\_site} & native & Double occupancy per site \\
\bottomrule
\end{longtable}

\subsection{RDM fields (features)}
\begin{longtable}{@{}p{2.4cm}p{1.6cm}p{9.3cm}@{}}
\toprule
Field & Dataset & Description \\ \midrule \endhead
\code{sites}       & both & Site cluster of the RDM \\
\code{n\_body}     & both & Cluster size $n$; RDM is $2^n\times2^n$ \\
\code{rho\_real}   & both & Real part of $\rho$ (npz array) \\
\code{rho\_imag}   & both & Imag part of $\rho$ (npz array) \\
\code{eigvals}     & both & Eigenvalues of $\rho$ (npz array) \\
\code{rdm\_eigvals} & both & Eigenvalues of $\rho$ (manifest, stringified) \\
\code{purity}      & both & $\mathrm{Tr}(\rho^2)$ \\
\code{is\_pure}    & both & purity $>1-10^{-6}$ \\
\code{entropy}     & both & von Neumann entropy $-\mathrm{Tr}(\rho\log\rho)$ \\
\code{tr\_rho2}    & both & $\mathrm{Tr}(\rho^2)$ \\
\code{tr\_rho3}    & both & $\mathrm{Tr}(\rho^3)$ \\
\code{rank}        & both & Numerical rank of $\rho$ \\
\code{raw\_trace}  & native & Trace of $\rho$ before normalisation \\
\bottomrule
\end{longtable}

\subsection{GME label fields}
\begin{longtable}{@{}p{2.9cm}p{1.6cm}p{8.8cm}@{}}
\toprule
Field & Dataset & Description \\ \midrule \endhead
\code{C\_target}            & both   & \textbf{Training label} $=$ \CGME{} (GCG for the gcg dataset) \\
\code{C\_gcg}               & gcg    & GCG convex-roof \CGME{} \\
\code{C\_pure}              & both   & Pure-state formula (exact roof if $\rho$ pure) \\
\code{lower\_bound\_super}  & both   & Supervisor analytical lower bound \\
\code{C\_UB}                & both   & Upper bound on \CGME{} \\
\code{C\_stiefel\_best}     & native & Best \CGME{} over Stiefel restarts \\
\code{C\_stiefel\_mean}     & native & Mean over restarts \\
\code{C\_stiefel\_std}      & native & Std over restarts (label noise) \\
\code{C\_stiefel\_n\_valid} & native & Number of valid restarts \\
\code{C\_stiefel\_total\_time} & native & Total Stiefel wall time \\
\code{C\_stiefel\_all}      & native & \CGME{} from every restart (npz) \\
\code{target\_source}       & native & \code{"pure\_formula"} or \code{"stiefel"} \\
\code{target\_in\_bracket}  & both   & $\text{LB}\le C_{\text{target}}\le\text{UB}$ \\
\code{status}               & gcg    & Solver status string \\
\code{solver}               & gcg    & Always \code{"gcg"} \\
\bottomrule
\end{longtable}

\subsection{Thermal fields (gcg)}
\begin{longtable}{@{}p{2.4cm}p{1.6cm}p{9.3cm}@{}}
\toprule
Field & Dataset & Description \\ \midrule \endhead
\code{T}       & gcg & Temperature (0 = ground state; $>0$ = thermal mixed) \\
\code{thermal} & gcg & True if the RDM is a thermal (mixed) state \\
\bottomrule
\end{longtable}

\subsection{Provenance / bookkeeping}
\begin{longtable}{@{}p{2.4cm}p{1.6cm}p{9.3cm}@{}}
\toprule
Field & Dataset & Description \\ \midrule \endhead
\code{config\_hash} & both   & Unique 12-char hash of the physics point (dedup key) \\
\code{sweep\_tag}   & both   & Origin: \code{vary\_U}/\code{vary\_t} or \code{ingested\_gme}/\code{surface}/\code{bank} \\
\code{source}       & gcg    & \code{gme\_runs} / \code{surface} / \code{bank} \\
\code{config\_json} & both   & Full config dict serialised (npz) \\
\code{label\_time\_s} & native & Wall time to compute the labels \\
\code{seed\_base}   & native & RNG seed (reproducibility) \\
\code{npz\_path}    & both   & Relative path to the \code{.npz} record \\
\bottomrule
\end{longtable}

% ===========================================================================
\section{Auditing the dataset}\label{sec:audit}
% ===========================================================================
\code{ml\_audit\_dataset.py} runs a battery of checks and prints a
\textsc{pass}/\textsc{warn}/\textsc{fail} verdict; it is meant to be re-run
before \emph{every} training run.
\begin{enumerate}[nosep]
  \item \textbf{Label completeness (NaN scan).} Counts NaNs in each label column;
        \textsc{fail} if a large fraction of \code{C\_target} is missing.
  \item \textbf{Bracket consistency.} Verifies
        $\text{lower\_bound\_super}\le C_{\text{target}}\le C_{\text{UB}}$ per record;
        flags violations (physical impossibility signals a solver bug).
  \item \textbf{\CGME{} distribution.} Histogram / quantiles of \code{C\_target};
        warns on degenerate (all-equal or all-zero) or too-narrow spreads that
        would make regression trivial or ill-posed.
  \item \textbf{Coverage.} Balance across $(L, n_{\text{body}}, U/t\text{ bin})$
        cells; warns if some physics regimes are under-represented.
  \item \textbf{Label noise.} When multiple solver restarts are recorded, reports
        the spread as an estimate of intrinsic label noise.
  \item \textbf{Pure vs mixed split.} Fraction of pure ($\text{purity}\approx1$)
        vs mixed records --- mixed states are the hard, informative cases.
\end{enumerate}

% ===========================================================================
\section{Commands}\label{sec:commands}
% ===========================================================================
Run from the repository root (the directory containing \code{main.py}).

\paragraph{Stage 1 --- build the RDM banks (general, cached in \code{results/}).}
\begin{Verbatim}[fontsize=\small]
for L in 2 4 6 8; do
  for t in 0.25 0.5 0.75 1.0 1.5 2.0; do
    for U in 2.0 4.0 8.0 10.0; do
      python main.py --build-bank --L $L --Ne $L --t $t --U $U --min-n 2 \
        --workers 8 --use-pkl-cache
    done
  done
done
# add --full-hamiltonian for the all-Sz sector; even L only for Sz=0 half-filling
\end{Verbatim}

\paragraph{Stage 2 --- label every bank with GCG (produces the dataset).}
\begin{Verbatim}[fontsize=\small]
python gme_qml/ml_ingest_pkls.py --source banks --bank-T 0 0.5 1 2 \
    --n-restarts 10 --max-iter 200 --k-max-extra 4
# --source all  also folds in existing gme_runs/ and surface PKLs
\end{Verbatim}

\paragraph{Stage 3 --- audit before training.}
\begin{Verbatim}[fontsize=\small]
python gme_qml/ml_audit_dataset.py --outdir gme_qml/data_for_ML/ingested --plots
\end{Verbatim}

\paragraph{Stage 4 --- train the baselines.}
\begin{Verbatim}[fontsize=\small]
python gme_qml/train_baselines.py --data gme_qml/data_for_ML/ingested
\end{Verbatim}

\paragraph{Cost note.} Stage~1 (banks) is cheap; Stage~2 costs one GCG solve per
record and is the dominant work. \code{-{}-bank-T 0} (ground state only) is the
cheapest; each additional temperature multiplies the number of solves by the
number of body orders. Memory grows combinatorially with $L$: the full-spectrum
bank is feasible up to $L\approx8$ on a large-RAM node; $L=10$ is infeasible.

% ===========================================================================
\section{Memory utilisation}\label{sec:memory}
% ===========================================================================
All density matrices are stored in \code{complex128}: \textbf{16 bytes per
matrix element} ($8$-byte real $+\,8$-byte imaginary). Every memory figure below
follows directly from this.

\subsection{A single $n$-body RDM (depends only on $n$)}
An $n$-body RDM is a $2^{n}\times2^{n}$ matrix, i.e. $4^{n}$ complex elements:
\[
  M_{\mathrm{RDM}}(n) \;=\; \underbrace{2^{n}\times 2^{n}}_{4^{n}\ \text{elements}}
  \times\, 16\ \text{bytes} \;=\; 4^{n}\cdot 16\ \text{bytes},
\]
independent of $L$. It \emph{quadruples} for every $+1$ in $n$.
\begin{center}
\begin{tabular}{@{}rrrr@{}}
\toprule
$n$ & matrix $2^n\!\times\!2^n$ & elements $4^{n}$ & memory \\ \midrule
2  & $4\times4$        & 16          & 256 B  \\
3  & $8\times8$        & 64          & 1 KB   \\
4  & $16\times16$      & 256         & 4 KB   \\
5  & $32\times32$      & 1\,024      & 16 KB  \\
6  & $64\times64$      & 4\,096      & 64 KB  \\
7  & $128\times128$    & 16\,384     & 256 KB \\
8  & $256\times256$    & 65\,536     & 1 MB   \\
9  & $512\times512$    & 262\,144    & 4 MB   \\
10 & $1024\times1024$  & 1\,048\,576 & 16 MB  \\
\bottomrule
\end{tabular}
\end{center}

\subsection{Full RDM bank and diagonalisation (depend on $L$)}
The bank holds one RDM \emph{per eigenstate} for every $n=2,\dots,L$, so it
scales with the Hilbert dimension $d(L)$. For the $S_z{=}0$ sector at half-filling
($N_e=L$, \emph{even} $L$),
\[
  d(L)=\binom{L}{L/2}^{2},\qquad
  M_{\mathrm{bank}}(L)=d(L)\sum_{n=2}^{L}4^{n}\cdot 16\ \text{bytes},
\]
and dense diagonalisation additionally needs the eigenvector matrix,
$M_{\mathrm{eigh}}(L)\approx 2\,d(L)^{2}\cdot 16$ bytes. Both grow
combinatorially; the $n{=}L$ term dominates the bank.
\begin{center}
\begin{tabular}{@{}rrrrrrl@{}}
\toprule
$L$ & $N_e$ & $d(L)$ & bank$[n{=}L]$ & full bank & dense eigh & feasibility \\ \midrule
2  & 2  & 4       & 1 KB    & 1 KB    & 512 B  & trivial \\
4  & 4  & 36      & 144 KB  & 189 KB  & 41 KB  & trivial \\
6  & 6  & 400     & 25 MB   & 33 MB   & 4.9 MB & laptop \\
8  & 8  & 4\,900  & 4.8 GB  & 6.4 GB  & 733 MB & large-RAM node \\
10 & 10 & 63\,504 & 992 GB  & 1.3 TB  & 120 GB & \textbf{infeasible} \\
\bottomrule
\end{tabular}
\end{center}
\emph{Odd $L$ at half-filling has no $S_z{=}0$ sector} ($N_e$ is odd, so
$N_\uparrow=N_\downarrow=N_e/2$ is non-integer); use the full Hamiltonian (all
$S_z$, dimension $\binom{2L}{N_e}$: $20, 252, 3432, 48\,620$ for $L=3,5,7,9$),
which is larger than the even-$L$ $S_z{=}0$ figures above. The single dominant
cost is always the $n{=}L$ bank --- e.g.\ $\text{bank}[n{=}10]=63\,504\times16$ MB
$=992$ GB on its own. On a $\sim$1.5 TB node the practical ceiling is
$L\approx8$; per-worker eigenvector duplication (avoided here by fork
copy-on-write) would otherwise multiply the eigenvector cost by the number of
workers.

\subsection{How the memory grows (scaling)}\label{sec:scaling}
There are two independent growth knobs, $n$ and $L$, and they compound.

\paragraph{Growing $n$: $\times4$ per step.}
An $n$-body RDM is a square matrix of side $2^{n}$. Adding one site to the
cluster \emph{doubles} the side ($2^{n}\!\to\!2^{n+1}$) and therefore
\emph{quadruples} the number of elements ($4^{n}\!\to\!4^{n+1}$). Memory
$=$ elements $\times16$ bytes, so it grows $\times4$ at every step:
\[
  \underbrace{256\,\text{B}}_{n=2}\!\to 1\,\text{KB}\!\to 4\,\text{KB}\!\to
  16\,\text{KB}\!\to 64\,\text{KB}\!\to 256\,\text{KB}\!\to 1\,\text{MB}\!\to
  4\,\text{MB}\!\to\underbrace{16\,\text{MB}}_{n=10},
  \qquad \text{each arrow } \times4 .
\]
This is pure exponential in $n$: $M_{\mathrm{RDM}}(n)\propto 4^{n}$.

\paragraph{Growing $L$: two exponentials multiply.}
The full bank is much steeper because two effects compound:
\begin{enumerate}[nosep]
  \item \textbf{More eigenstates.} The bank stores one RDM \emph{per eigenstate},
        and the number of eigenstates is
        $\;d(L)=\binom{L}{L/2}^{2}\approx 4^{L}/L$:
        \[
          d:\quad 4 \to 36 \to 400 \to 4\,900 \to 63\,504
          \qquad(\times9,\ \times11,\ \times12,\ \times13\ \text{per }{+}2).
        \]
  \item \textbf{Bigger RDMs.} Each state's largest RDM ($n{=}L$) is itself
        $4^{L}\cdot16$ bytes, growing $\times16$ per ${+}2$ in $L$.
\end{enumerate}
Since $\text{bank}\approx(\text{number of states})\times(\text{size of the largest RDM})$,
\[
  M_{\mathrm{bank}}(L)\;\approx\; d(L)\,\cdot\,4^{L}\cdot 16
  \;\approx\;\frac{4^{L}}{L}\cdot 4^{L}\;\sim\;\frac{16^{\,L}}{L}.
\]
So each ${+}2$ in $L$ multiplies memory by roughly
$\underbrace{\times12}_{\text{more states}}\times\underbrace{\times16}_{\text{bigger RDMs}}\approx\times200$,
exactly the observed ladder:
\[
  \underbrace{1\,\text{KB}}_{L=2}\!\to 189\,\text{KB}\!\to 33\,\text{MB}\!\to
  6.4\,\text{GB}\!\to\underbrace{1.3\,\text{TB}}_{L=10},
  \qquad \text{each arrow } \approx\times200 .
\]

\paragraph{Takeaway.}
\begin{center}
\begin{tabular}{@{}lll@{}}
\toprule
Knob & Scaling & Growth rate \\ \midrule
$n$ (a single RDM)   & $\propto 4^{n}$  & $\times4$ per step \\
$L$ (the full bank)  & $\propto 16^{L}/L$ & $\times{\sim}14$ per step ($\times200$ per ${+}2$) \\
\bottomrule
\end{tabular}
\end{center}
$L$ is the killer: it is the \emph{product} of two exponentials --- more states
to store \emph{and} each state's RDM exponentially larger. This is why a single
$n{=}10$ RDM is only $16$ MB, yet the whole $L{=}10$ bank is $1.3$ TB: the
$n{=}10$ slice alone is $16\,\text{MB}\times 63\,504\ \text{states}\approx 1$ TB.

% ===========================================================================
\section{Current status: what is already done}\label{sec:status}
% ===========================================================================
This section summarises, for the supervisor, exactly what data exists and what
has been trained on it so far.

\subsection{Dataset inventory}
\begin{center}
\begin{tabular}{@{}lrl@{}}
\toprule
Source & Records & Coverage \\
\midrule
Native manifest (\code{ml\_generate\_ml\_dataset.py}) & 2{,}610 &
  $L=3$--$6$, uniform ring, vary-$U$ + vary-$t$, ground + excited states,
  $n=2..L$ \\
Ingested from PKL caches (\code{ml\_ingest\_pkls.py}) & 111 &
  GCG-labelled records mined from existing RDM banks / surface / thermal PKLs
  (free data, no recomputation) \\
\bottomrule
\end{tabular}
\end{center}
All records are Sz$=0$ sector; ground- and excited-state (T$=0$) and thermal
(mixed, $T>0$) states are both represented. The audit (\code{ml\_audit\_dataset.py})
passes on completeness, RDM validity and coverage, with one open flag: 57/2610
(2.2\%) native records show a bracket violation
($\text{LB}\le C_{\text{target}}\le\text{UB}$ broken), traced to the label
solver (Section~\ref{sec:gaps}).

\subsection{Trained baselines (Stage 4, \code{train\_baselines.py})}
Evaluated with a \textbf{group split by Hamiltonian} (all records from the same
$(L,t,U,\text{state})$ stay in one fold, preventing leakage from near-duplicate
site-clusters): 1{,}992 training / 660 held-out test records, 103/35 groups.

\paragraph{Task A --- Detection ($\CGME{}>\varepsilon=0.01$).}
\begin{center}
\begin{tabular}{@{}lccccc@{}}
\toprule
Model & Accuracy & ROC-AUC & F1 & Hard-band acc.\ ($n{=}181$) \\
\midrule
Logistic regression & 0.952 & 0.960 & 0.969 & 0.823 \\
Random forest        & 0.955 & 0.973 & 0.971 & 0.834 \\
\textbf{HistGBM (best)} & \textbf{0.962} & 0.972 & \textbf{0.976} & \textbf{0.862} \\
\bottomrule
\end{tabular}
\end{center}
The \emph{hard band} (records within $0.05$ of the threshold $\varepsilon$) is
the honest metric: overall accuracy is inflated by the many unambiguous
(strongly entangled or clearly separable) records. An 86\% hard-band accuracy
means roughly 1 in 7 boundary cases is still misclassified --- consistent with
the literature, where near-threshold detection is the known weak point of all
ML entanglement classifiers.

\paragraph{Task B --- Estimation (regressing the value of \CGME{}).}
\begin{center}
\begin{tabular}{@{}lccc@{}}
\toprule
Model & MAE & RMSE & $R^2$ \\
\midrule
Physics baseline (supervisor's analytic LB) & 0.166 & 0.273 & 0.443 \\
Ridge                    & 0.150 & 0.208 & 0.676 \\
Random forest            & 0.049 & 0.096 & 0.932 \\
\textbf{HistGBM (best)}  & \textbf{0.042} & 0.094 & \textbf{0.934} \\
\bottomrule
\end{tabular}
\end{center}
The best model (HistGBM) reduces the mean error to a quarter of the physics
baseline's ($0.042$ vs.\ $0.166$). Broken down by body order $n$ (HistGBM):
\begin{center}
\begin{tabular}{@{}cccc@{}}
\toprule
$n$ & MAE & $R^2$ & note \\
\midrule
2 & 0.014 & 0.994 & concurrence-equivalent; near-exact \\
\textbf{3} & \textbf{0.086} & \textbf{0.698} & \textbf{weakest --- the minimal genuinely-multipartite order} \\
4 & 0.017 & 0.990 & near-exact \\
\bottomrule
\end{tabular}
\end{center}
The degradation at $n=3$ is the single most important finding so far: it is
precisely the smallest cluster size at which \emph{genuine} multipartite
(as opposed to pairwise) entanglement exists, and it is where the RDM spans the
richest mix of separable / biseparable / genuinely-entangled structure. This is
the priority target for the next modelling iteration (better features and/or
more $n=3$ training data).

% ===========================================================================
\section{What data is still required}\label{sec:gaps}
% ===========================================================================
Ordered by priority.
\begin{enumerate}[nosep]
  \item \textbf{Relabel with GCG.} The 2{,}610 native records were labelled by
        \code{convex\_roof\_stiefel}, which we have since verified is less
        accurate than GCG (the current default solver everywhere else in the
        project). The label noise this introduces is the most likely cause of
        the 2.2\% bracket violations. \emph{Action:} add a
        \code{-{}-label-solver gcg} mode to \code{ml\_generate\_ml\_dataset.py}
        (or fully rely on the GCG-labelled ingestion path) and relabel/retrain.
  \item \textbf{More separable ($\CGME{}\approx0$) examples.} Ground-state
        Hubbard RDMs are almost always entangled; the only source of
        near-zero labels is high-temperature thermal states. \emph{Action:}
        ingest denser, hotter $T$ grids from the thermal/surface PKLs to fix
        the class imbalance feeding the hard-band weakness above.
  \item \textbf{Full-Hamiltonian (all-$S_z$) records.} Every record so far is
        Sz$=0$. The sector tag already exists in the caches, so ingesting
        \code{-{}-full-hamiltonian} banks/runs requires no new physics code,
        only running the ingestion over those caches.
  \item \textbf{More $n=3$ coverage.} Given the finding above, additional
        $n=3$ records (more $U/t$ points, more thermal states, and/or the new
        molecular 3-orbital clusters) are the highest-value single addition to
        the dataset.
  \item \textbf{Molecular records (new source).} The molecular GME pipeline
        (Benzene CAS$(6,6)$, canonical and Pipek--Mezey localized bases, T$=0$
        and thermal) now produces GCG-labelled RDMs in the same PKL format.
        These are not yet ingested; adding them would let the dataset test
        transfer from the Hubbard ring to a real molecular system.
\end{enumerate}

% ===========================================================================
\section{Summary}
% ===========================================================================
The dataset couples a cached, task-independent \emph{RDM bank} (the feature
source) with the \emph{GCG convex roof} (the label). Banks are built once per
Hamiltonian in \code{results/}; records --- RDM $+$ derived features $+$ \CGME{}
--- are written to \code{gme\_qml/data\_for\_ML/ingested/} and audited before
training. Every stored field is enumerated in Section~\ref{sec:schema} and in
the machine-readable \code{data\_dictionary.csv}. Baselines are trained
(Section~\ref{sec:status}); the priority list for the next data-generation pass
is given in Section~\ref{sec:gaps}, headed by relabelling with GCG and
improving $n=3$ / class-balance coverage.

\end{document}
